\documentclass{UAmathtalk}

\author{Barry Minemyer}
\address{Commonwealth University-Bloomsburg}
\urladdr{https://facstaff.bloomu.edu/bminemyer/}
\title{Complex Hyperbolic Gromov–Thurston Metrics}
\date{Thursday, November 19, 2026}

\renewcommand*{\where}{Massry B008}


\begin{document}

\maketitle

\begin{abstract}
In 1987 Gromov and Thurston discovered the first Riemannian manifolds that are not homotopy equivalent to a hyperbolic manifold, but admit a Riemannian metric whose sectional curvatures are pinched arbitrarily close to~$-1$. The manifolds that they construct are branched covers of hyperbolic manifolds, and to construct the metric they perform a sort of “geometric surgery” about the ramification locus. In 2022 Stover and Toledo proved the existence of similar branched cover manifolds built out of complex hyperbolic manifolds, and via a result of Zheng these manifolds admit a negatively curved Kahler metric. In this talk we will discuss how to construct a (not Kahler) Riemannian metric on these Stover–Toledo manifolds which is arbitrarily close to being negatively $1/4$-pinched. This provides the first known examples of Kahler manifolds that are not homotopy equivalent to a complex hyperbolic manifold, but admit a Riemannian metric with this property.
\end{abstract}

\end{document}
